%% sections/09_mechanism_analysis.tex

\section{메커니즘 분석}
\label{sec:mechanism}

\subsection{왜 naive n-gram은 실 트레이스에서 붕괴하는가}

prompt-lookup n-gram 은 입력 내 n-gram 일치에만 의존한다.
실 서빙 trace 는 고유 식별자$\cdot$다국어 토큰$\cdot$비반복 구문이 지배적이어서
hit rate 가 극히 낮고, $\alpha K_{\mathrm{eff}}\ll\mathrm{overhead}$ 가 되어
검증 비용만 누적된다. Qwen-7B 에서 n-gram 은 corpus 무관하게 $\sim\!-90\%$ 로 붕괴한다.
따라서 naive n-gram 은 게이트 후보에서 제외된다.

\subsection{왜 suffix는 대부분의 모델에서 작동하는가}

SuffixDecoding 은 입력과 이전 생성 출력을 suffix tree 에 누적하여
반복 구조$\cdot$정형 구문(코드$\cdot$JSON)$\cdot$라이브러리 호출을 드래프트로 재사용한다.
그 결과 dense 표준 모델에서 $\alpha=0.6\!\sim\!0.93$ 의 높은 수락률을 달성하고
(Table~\ref{tbl:b200-latency}), $\alpha$ 가 높을수록 이득이 커지는 R/K 관계가
10개 모델에서 일관 관찰된다(Fig.~\ref{fig:alpha-gain}).

\subsection{왜 추론 특화 모델에서 $\alpha$가 무너지는가}
\label{subsec:mech-reasoning}

추론 특화 모델(R1 계열)은 답 이전에 긴 사고 사슬을 생성한다.
같은 instruction 에도 서로 다른 추론 경로를 취하므로 출력의 표층 반복성이 낮고,
suffix tree 에 축적된 과거 출력과의 prefix-match 율이 떨어진다.
실측에서 R1-Distill-Llama-70B 의 단일 corpus $\alpha$ 는 $0.27\!\sim\!0.38$ 로,
$1+\alpha K_{\mathrm{eff}}\approx 1.3$ --- 즉 spec 효과가 거의 소멸한다.
overhead 자체는 작아(wall ratio $\approx 1.1$) 손실 폭은 $-2\!\sim\!-15\%$ 로 작지만
부호는 일관되게 음(陰)이다.
(mix 조건만 net-positive 인 이유는 누적 cache 의 cross-corpus 부분 일치로
$\alpha$ 가 일부 회복되기 때문이다.)

\subsection{왜 MoE 추론(R1-671B)이 최악인가}
\label{subsec:mech-moe}

R1-671B 는 \emph{두 효과}를 결합한다.
(i) 추론 출력 다양성으로 $\alpha$ 가 낮고($\approx0.45$, \S\ref{subsec:mech-reasoning}),
(ii) MoE 는 토큰마다 expert 라우팅이 달라 draft 토큰열의 분포 예측을 더 어렵게 하며,
동시에 671B activated 경로의 \emph{verify forward 비용이 커서 overhead 가 크다}
(wall ratio $1.6\!\sim\!2.9$).
식~\ref{eq:rk-condition}의 두 항이 모두 불리한 방향이므로,
$\alpha\approx0.5$ 에서도 $-40\!\sim\!-64\%$ 로 크게 손해다(Table~\ref{tbl:b200-xl}).
이는 \emph{최적 방법이 크기가 아니라 모델 유형}임을 메커니즘 수준에서 설명한다.

\subsection{성공/실패가 자원에 남기는 상반된 흔적}
\label{subsec:mech-resource}

부호는 자원 균형으로도 드러난다(\S\ref{subsec:res-resource}).
\textbf{성공}하면 step 수가 줄어 GPU 가 un-saturate 되고($94.9\!\to\!78.7\%$),
줄어든 step 사이의 host-path 작업이 상대적으로 커진다 ---
시스템이 host(CPU)-bound 로 전환되며 회수 가능한 GPU 여유와 유휴 CPU 가 공존한다.
\textbf{실패}하면 거부되는 draft 의 낭비 연산으로 GPU 가 포화를 유지하므로
($98.2\!\to\!86.8\%$) host-bound 여유가 작다.
따라서 host-side reclamation(\S\ref{sec:direction})의 기회는
\emph{정확히 spec-decode 가 성공하는 모델$\cdot$조건}에 비례하여 존재한다.
이 인과를 정량 분해하는 것이 \S\ref{subsec:decisive} 의 결정적 실험이다.
